% !PRINTED 678
\weberpage{678}
%<<CONT
be proved for $s - 1$ and derive it in general by complete induction. To this end we arrange the determinant $L_s$ according to the elements of the last row, and write it thus:
\weq{2}{
L_s = g_1 l_{s,1} + g_2 l_{s,2} + \cdots + g_s l_{s,s},
}
in which $g_s = L_{s-1}$, which can therefore, by the assumption made, be taken to be different from zero. If we substitute these values of $g$ into the formula of Theorem 4., \S.\,185, assuming $g_{s+1} = 0, \ldots, g_\nu = 0$ while $g_s$ is not $= 0$, then, so long as $s < \nu$, the quantities $g_1, g_2, \ldots, g_\nu$ are certainly not all equal to one another, and it follows that $\varepsilon_s$ can be so chosen that $L_s$ is different from zero, as was to be proved.

This mode of reasoning cannot be extended to $s = \nu$, and the determinant $L_\nu$ must indeed always be zero, because for every unit $\sum\limits_{1,\nu}^{s} l_s$ vanishes.

We now wish, for the case that $s = \nu - 1$, to denote the determinant $L_s$ thus:
\weq{3}{
L_{\nu-1} = L(\varepsilon_1, \varepsilon_2, \ldots, \varepsilon_{\nu-1}),
}
and to call a system of units $\varepsilon_1, \varepsilon_2, \ldots, \varepsilon_{\nu-1}$ for which this determinant is different from zero a \gesperrt{system of independent units}.

The absolute value $L$ of the determinant $L(\varepsilon_1, \varepsilon_2, \ldots, \varepsilon_{\nu-1})$ is called the \gesperrt{regulator} of the system $\varepsilon_1, \varepsilon_2, \ldots, \varepsilon_{\nu-1}$\footnote{\gesperrt{Dirichlet-Dedekind}, Vorlesungen \"uber Zahlentheorie 4. Auflage, \S.\,183.}.

If in the determinant
\weq{4}{
\begin{vmatrix}
l_{1,1}, & l_{1,2}, & \ldots, & l_{1,\nu} \\
\cdots & \cdots & \cdots & \cdots \\
l_{\nu-1,1}, & l_{\nu-1,2}, & \ldots, & l_{\nu-1,\nu} \\
x_1, & x_2, & \ldots, & x_\nu
\end{vmatrix},
}
in which the $x_1, x_2, \ldots, x_\nu$ are arbitrary quantities, we add all the columns to the last, then, having regard to the relations $\sum\limits^{i} l_{s,i} = 0$
we obtain its value equal to
\[
\pm (x_1 + x_2 + \cdots + x_\nu) L.
\]
If therefore we let all the $x$ with the exception of one arbitrary
%CONT>>
